Quantum-inspired encodings (QIE) — amplitude, angle, and basis encoding —
have been proposed as fixed feature maps that import quantum-circuit geometry
into classical machine learning pipelines.
We conduct the first systematic benchmark of all three encodings across ten
diverse datasets, evaluating predictive performance, spectral geometry, kernel
alignment, and practical overhead under a strictly matched representational budget.
Across 30 encoding-dataset pairs, QIE achieves zero statistically significant
accuracy wins against strong classical baselines; 27 are significantly worse at
$\alpha{=}0.05$, and the remaining three are negligible near-ties
($|d|\leq 0.18$, $p\geq 0.70$; paired $t$-test and Wilcoxon signed-rank).
Rather than reporting a simple failure, we provide a mechanistic spectral
attribution for each encoding: amplitude encoding undergoes rank collapse under
$\ell_2$-normalisation ($\erank$ as low as 1.04, condition number $\kappa$
up to $5.7{\times}10^9$); angle encoding is geometrically equivalent to a
scaled linear map under logistic probing
($\mathrm{CKA}(\text{angle, raw linear})\geq 0.95$ on 7/10 datasets);
and basis encoding, while spectrally well-conditioned, imposes a binary Hamming
geometry misaligned with smooth discriminative manifolds.
The one structural condition under which QIE approaches competitive performance,
namely uniformly filled output space, is exemplified by amplitude encoding on the
high-rank noise dataset ($\erank{=}197.3$, $d{=}{-}0.03$, $p{=}0.96$) and is
precisely predicted by the spectral framework.
Our findings imply that the geometry imposed by quantum-circuit-inspired feature
maps does not, in isolation, confer a representational advantage over classical
alternatives, and that encoding overhead is not the limiting factor.
